\chapter*{Abstract}

Malaria remains a leading cause of morbidity and mortality among children under five in sub-Saharan Africa. Insecticide-treated nets (ITNs) are a primary malaria prevention intervention, but estimating their causal effect using observational data is challenging due to confounding by environmental, socioeconomic, and demographic factors.

This study estimates the causal effect of ITN use on malaria infection among children under five in Kenya using data from the 2015 and 2020 Kenya Malaria Indicator Surveys. The analytical sample comprises 6,771 children across 528 clusters. The study employs multiple causal inference methods: traditional logistic regression, Generalized Linear Mixed Models (GLMM), propensity score matching, inverse probability of treatment weighting, doubly robust estimation, and Double Machine Learning. Unlike traditional approaches, propensity scores are estimated using GLMM to account for the hierarchical structure of the data, which exhibits an intra-cluster correlation of 41-47\%.

The results reveal substantial confounding. Unadjusted models show a harmful association, while fully adjusted models reveal a protective effect. GLMM-based causal methods produce statistically significant protective effects ranging from 45-55\% reduction in malaria odds. Double Machine Learning provides more conservative estimates, with odds ratios between 0.961 and 0.981 representing a 2-4\% reduction at the population level. Heterogeneity analysis shows that ITN effectiveness varies by subgroup, with strongest effects observed in the poorer wealth quintile, urban areas, older children aged 36-59 months, and high-rainfall regions.

The study concludes that ITN use has a protective causal effect on malaria infection among children under five in Kenya. The triangulation of evidence across six methods strengthens confidence in this conclusion. The distinction between cluster-specific effects (45-55\% reduction) and population-average effects (2-4\% reduction) highlights the importance of understanding which estimand is relevant for different policy questions. Methodologically, the study demonstrates that accounting for clustering in propensity score estimation is essential for valid causal inference. Policy implications include continued investment in ITN distribution programs with targeted strategies for high-risk populations and complementary interventions for the poorest households.

\vspace{0.5cm}

\textbf{Keywords}: Malaria, Insecticide-treated nets, Causal inference, Double Machine Learning, Propensity score matching, Kenya, Children under five


\vfill
\section*{Declaration}
I, the undersigned, hereby declare that the work contained in this essay is my original work, and that any work done by others or by myself previously has been acknowledged and referenced accordingly.

% Scan your signature into a small picture called 'signature.png' and insert it
% above your name and the date:
\includegraphics[height=4cm]{Figures/signature} \hrule
% Your name must be in English Capitalisation with no comma, and the Family name comes last. 
% Do note the date below. It is called the "deadline".
Nicholas Mwanza, 30 May 2026.
